% Чуть компактнее вертикальные интервалы, чтобы завершить билет на одной странице.
\begingroup
\setlength{\parskip}{2pt}
\tcbset{before skip=4pt,after skip=4pt,top=5pt,bottom=5pt}
\source{Лекции, с. 32–37; Карпов, §6.1, §6.3; детали индукции восстановлены составителем}
\begin{definition}[Укладка на поверхности и грани]
Пусть \(M\) — поверхность. Укладка графа \(G\) на \(M\) состоит
из инъекции \(\varphi_V:V\to M\) и простых дуг
\(\varphi_e:[0,1]\to M\), \(e=uv\), с концами
\(\varphi_e(0)=\varphi_V(u)\), \(\varphi_e(1)=\varphi_V(v)\).
Внутренность дуги не содержит образов вершин, а две дуги могут встречаться
только в образах общих концов соответствующих рёбер.
Планарный граф допускает укладку на плоскости; плоский граф — граф
с уже выбранной укладкой. Грани — компоненты дополнения к рисунку;
на плоскости ровно одна грань неограниченна и называется внешней.
Число граней обозначается \(f\). Можно работать с кусочно-линейными дугами.
\end{definition}
\begin{lemma}[Удаление рёбер на рисунке]
Ребро цикла граничит с двумя различными гранями; его удаление объединяет
эти грани и уменьшает \(f\) на 1. Висячее ребро не разделяет грань;
его удаление вместе с листом сохраняет число граней.
\end{lemma}
\begin{proof}
Простой цикл через ребро — кривая Жордана. Две стороны внутренней
точки ребра лежат по разные стороны этой кривой и потому не могут
принадлежать одной грани. Малый открытый коридор на месте удалённого
ребра соединяет именно эти две области; других изменений дополнения нет.
У висячего ребра обе стороны соединяются обходом свободного конца
в малой окрестности, не содержащей других рёбер. Оно лишь образует
разрез внутри одной области, а не отделяет новую область.
\end{proof}
\begin{theorem}[Формула Эйлера]
Для плоского графа с \(c\) компонентами
\[
 \boxed{n-m+f=1+c.}
\]
В частности, для связного непустого графа \(n-m+f=2\).
\end{theorem}
\begin{proof}
Сначала пусть граф связен. Строим его рисунок из одной вершины.
По билету 6 есть остовное дерево; укореним его и добавим остальные
вершины по одной вместе с ребром к уже построенному родителю.
Числа \(n,m\) возрастают на 1, число граней не меняется: висячее ребро
грань не разделяет. Затем добавляем остальные рёбра в их исходных дугах.
Концы очередного ребра уже соединены остовом, поэтому новое ребро
замыкает простой цикл и делит одну грань на две по теореме Жордана.
Числа \(m,f\) возрастают на 1. Оба шага сохраняют \(n-m+f\).
В базе \(n=1,m=0,f=1\), значит, для связного графа \(n-m+f=2\).

\textit{Обобщение составителя.} Для несвязного графа можно соединить разные компоненты рёбрами без
пересечений: берём дугу между рисунками двух компонент, и участок между
последовательными касаниями разных компонент, внутренность которого
не встречает граф. Его концы при необходимости доводим вдоль границ
грани до вершин. Новое ребро — мост и не создаёт грани: его стороны
соединяются обходом одной из прежде раздельных компонент.
После \(c-1\) добавлений получим связный граф с теми же \(n,f\)
и \(m+c-1\) рёбрами. Отсюда \(n-(m+c-1)+f=2\).
Для пустого графа \(n=m=c=0,f=1\), равенство также верно.
\end{proof}
\begin{definition}[Подразбиение и гомеоморфные графы]
Элементарные операции: заменить ребро \(uv\) на \(ut,tv\) с новой
вершиной \(t\); обратно, удалить вершину степени 2, склеив две её дуги
в одно ребро (при необходимости разрешая промежуточные кратности).
Графы гомеоморфны, если связаны последовательностью таких операций
и изоморфизмов. Эквивалентно, они имеют изоморфные подразбиения.
Подразбиение — замена рёбер простыми путями с новыми внутренними
вершинами; у разных рёбер эти вершины различны. Это комбинаторное понятие;
гомеоморфизм плоскости — непрерывная биекция с непрерывной обратной.
\end{definition}
\begin{theorem}[Инвариантность планарности]
Граф планарен тогда и только тогда, когда планарно любое его подразбиение.
Следовательно, гомеоморфные графы одновременно планарны.
\end{theorem}
\begin{proof}
При подразделении отмечаем новые вершины на дугах рёбер. Обратно,
склеиваем последовательные дуги каждого заменяющего пути в одну простую
дугу; других пересечений при этом не возникает. Гомеоморфные графы
сравниваются через общее подразбиение.
\end{proof}

\endgroup
